\documentclass[11pt, a4paper, english, twoside, notitlepage]{report}

\usepackage{iipr}
\fancyhead[RO, LE]{ }
\fancyhead[LO]{\nouppercase \leftmark}	% Heading on odd pages: Chapter name on the left
\fancyhead[RE]{\nouppercase \rightmark} % Heading on even pages: Section on the right
\fancyfoot[LE,RO]{\thepage} % Roman numbering, foot left on even, foot right on odd pages
\fancyfoot[CE]{Doble Grado en Matem\'aticas e Ingenier\'ia Inform\'atica} %Foot: Center, even pages
\fancyfoot[CO]{Universidad Complutense de Madrid (UCM)} %Foot: Center, odd pages
\renewcommand{\footrulewidth}{0.4pt} % Rule on top of foot
\renewcommand{\headrulewidth}{0.4pt} % Rule below headings
\pagestyle{fancy}

% ________Pruebas________
\usepackage{lipsum}
% _______________________

\begin{document}

\title{\huge{\textsc{Complexity analysis of polynomial algorithms}} \\
	\protect\includegraphics[scale=0.2]{ucm.pdf}}
\author{Ignacio Iker Prado Rujas \\
	Universidad Complutense de Madrid (UCM) \\
	Doble Grado en Matem\'aticas e Ingenier\'ia Inform\'atica}
\date{\today}
\maketitle
\begin{center}
	Tutors: \\
	Jos\'e F. Fernando \textit{\&} J.M. Gamboa \\
	\
\end{center}

\thispagestyle{empty}	
	
\begin{abstract}
	This work is about three different proofs of the same fact and a computational comparison between them, looking for the best one.

	Let $R$ be a real closed field and $n\geq 2$. We prove that: (1) for every finite subset $F$ of $R^n$,
	the semialgebraic set $R^n\setminus F$ is a polynomial image of $R^n$; and (2) for any
	independent linear forms $l_1,\ldots,l_r$ of $R^n$, the semialgebraic set
	$\{l_1>0,\ldots,l_r>0\}\subset R^n$ is a polynomial image of $R^n$.
	
	The key proof here is that $Q = \{x > 0, y > 0\}$ is a polynomial image of $\R^2$. This assert is proved in three different ways: a first approach using real algebraic geometry; a second and shorter one, using the composition of 3 rather simple maps; and a third one that applies topology with no computer computations.
	
	(...)
\end{abstract}

\newpage
\mbox{}
\thispagestyle{empty}    % White page
\newpage

\begin{chapter}{Polynomial images of $R^n$}
\pagenumbering{arabic}    % Numbering from 1

\begin{section}{Introduction}
	\begin{definition}\label{polyMap}
	Let $R$ be a real closed field (see definition \ref{realCField}) and $m, n \in \N_{>0}$. A map $f = (f_1, \dots, f_n): R^m \longrightarrow R^n$ is said to be polynomial if $f_i \in R[x_1, \dots, x_m], i = 1,\dots, n$. 
	\end{definition}
	
	A very famous theorem by Tarski and Seidenberg states:
	\begin{theorem}[Tarski-Seidenberg]\label{tarskiSeidenberg}
	The image of any polynomial map $f: R^m \longrightarrow R^n$ is a semialgebraic subset (see definition \ref{semialgSet}) of $R^n$.
	\end{theorem}
	
	In this work we are studying sort of a converse of this statement. In an \emph{Oberwolfach} week, J.M. Gamboa proposed to characterize the semialgebraic subsets of $R^n$ that are polynomial images of $R^m$.
	%Comentar lo de los abiertos en relación con la conjetura Jacobiana?
	\begin{notation}
	We need to mention to which topology we refer when we talk about closures, boundaries, etc. More specifically, the \textbf{exterior boundary} of a set $S$ is $\delta S := \bar{S} \backslash S$, with $\bar{S}$ being the \textbf{closure} of $S$ in $R^n$ in the usual topology. $\bar{S}^{\text{zar}}$ is the closure of $S$ with respect to the Zariski topology (see definition \ref{zariski}). $A\subset R^n$ is \textbf{irreducible} if its Zariski closure $\bar{A}^{\text{zar}}$ is an irreducible algebraic set.
	\end{notation}
	
	\begin{subsection}{Neccesary conditions and examples}
	To begin working on this idea, we provide some neccesary conditions for a set $S\subset R^n$ to be polynomial image of $R^m$. 
	
	It is trivial that for $m = n = 1$ (so $f: R \rightarrow R$), the images of polynomial maps are either a set of one point or singletons (if the map is constant), or unbounded closed intervals (think of $f(x) = x^2$), or the whole $R$ (think of $f(x) = x$).
	
	In the general case, by theorem \ref{tarskiSeidenberg}, .
	
	\end{subsection}
	
\end{section}

\end{chapter}


\appendix
\begin{chapter}{Auxiliar definitions and results}

\begin{definition}\label{realCField}
A \textbf{real closed field} is a field $R$ that has the $1\textsuperscript{st}$ order properties as the field of real numbers $\R$.
\end{definition}

\begin{definition}\label{semialgSet}
A \textbf{semialgebraic set} is a subset $S \subset R^n$ (for some real closed field $R$) defined by a finite sequence of polynomial equations of the form:
\begin{align*}
	P_1(x_1, \dots&, x_n) = 0\\
	\vdots &\\
	P_r(x_1, \dots&, x_n) = 0 \\
	Q_1(x_1, \dots&, x_n) > 0 \\
	\vdots &\\
	Q_l(x_1, \dots&, x_n) > 0 \\
\end{align*}
. A \textbf{semialgebraic map} is a map that has semialgebraic graph. Moreover, the finite union, intersection and complement of semialgebraic sets is still a semialgebraic set.
\end{definition}

\begin{definition}[Zariski topology]\label{zariski}
It is a topology on algebraic varieties whose closed sets are the algebraic subsets of the variety. Its sets are defined as the set of solutions of a system of polynomial equations over a field $R$. In this topology, when we talk about the irreducibility of a element, we mean that it is not the union of two smaller sets that are closed under the Zariski topology.
\end{definition}

\end{chapter}


\begin{thebibliography}{10}

%\bibitem{...}

\end{thebibliography}

\end{document}